% =====================================================================
% Model definition -- first piece of the Methods section (DRAFT)
% Needs: \usepackage{amsmath,amssymb,booktabs}
% Citation keys used: johnson2014trophic, klaise2016neurons, degroot1974reaching
% (BibTeX entries at the bottom of this file, commented out)
% Lines starting with "% TODO" or "% CHECK" are notes to self; they do not render.
% =====================================================================

\subsection{Network generation}
\label{sec:generation}

A group is a directed graph $G=(V,E)$ on $N$ nodes, where $(i,j)\in E$ means
that $i$ has authority over $j$ ($i$ is senior, $j$ junior). $A$ is the
adjacency matrix, $k_i=\sum_j A_{ij}$ the out-degree of $i$, and $[P]$ the
Iverson bracket. This is the transpose of the ecological convention in which
the generalized preferential preying model (GPPM) and trophic coherence were
defined~\cite{johnson2014trophic,klaise2016neurons}, where $(i,j)$ means
$i$ is eaten by $j$. ``Predator'' here corresponds to ``senior'' and ``prey''
to ``junior''. Section~\ref{sec:equivalence} shows the two descriptions define
the same model.

The generator takes size $N$, basal fraction $b$, density $c$ and normalized
kernel width $\hat T$, and sets $B=\max(1,\operatorname{round}(bN))$ and
$L=\operatorname{round}(cN)$.

\paragraph{Tree phase.}
Nodes $0,\dots,B-1$ are basal, with level $\ell_i=1$. Nodes $n=B,\dots,N-1$
are added in order. Each picks a node already present, uniformly at random,
and becomes its superior:
\begin{equation}
  j_n\sim\mathrm{Unif}\{0,\dots,n-1\},\qquad
  E\leftarrow E\cup\{(n,j_n)\},\qquad
  \ell_n=\ell_{j_n}+1 .
  \label{eq:tree}
\end{equation}
This gives a forest of $B$ trees rooted at the basal nodes, with $N-B$ edges
$E_{\mathrm{tree}}$. Levels $\ell$ are never updated after this phase. Basal
nodes have $k_i=0$ and all others $k_i\ge1$, so $\ell_i=1\iff k_i=0$. Since
newcomers sit above the node they pick, later arrivals tend to be more senior.

\paragraph{Excess-edge phase.}
Let $D(G)=\max_i\ell_i$ be the height of the tree just built, and set
$T=\hat T\,D(G)$ (Section~\ref{sec:that}). A further $L-(N-B)$ distinct edges
are drawn from
\begin{equation}
  \mathcal C=\bigl\{(i,j):\ \ell_i\ge 2,\ i\ne j,\ (i,j)\notin E_{\mathrm{tree}}\bigr\}
  \label{eq:candidates}
\end{equation}
with probability proportional to a Gaussian kernel on the level gap, peaked at
a gap of one:
\begin{equation}
  \Pr[(i,j)]\ \propto\ w_{ij}=\exp\!\left[-\frac{(\ell_i-\ell_j-1)^2}{2T^2}\right].
  \label{eq:kernel}
\end{equation}
Basal nodes are excluded as sources but not as targets. A candidate may
reverse an existing edge, so reciprocated ties can form between non-basal
nodes. As $T\to0$ every excess edge spans one level; as $T\to\infty$ edges are
uniform over $\mathcal C$. Since $w_{ij}$ depends only on $(\ell_i,\ell_j)$, we
draw a level pair with probability proportional to its total weight and then a
candidate uniformly within it, with replacement, rejecting duplicates. This
successive sampling closely approximates weighted sampling without replacement
because $L-(N-B)\ll|\mathcal C|$.

\paragraph{Cleanup and density.}
We keep the largest weakly connected component, which contained at least 99\%
of nodes in every realization (mean above 99.99\%); $N$ denotes nominal size
throughout. The density is fixed across all sizes, so mean degree does not
change with $N$:
\begin{equation}
  c=(1-b)+\mu\,\frac{B_{\max}\ln B_{\max}}{N_{\max}},\qquad
  B_{\max}=\operatorname{round}(bN_{\max}).
  \label{eq:density}
\end{equation}
The first term counts tree edges per node. The second is a coupon-collector
estimate of the excess edges per node needed to join all $B$ trees, inflated
by a safety factor $\mu=9$. With $b=0.1$ and $N_{\max}=10^4$, $c\approx7.12$.

\subsection{Equivalence with the GPPM}
\label{sec:equivalence}

Our implementation differs from the published GPPM~\cite{klaise2016neurons}
in three places, each a reversal of direction (Table~\ref{tab:reversals}).
All three follow from the change of convention above. Let
$R(E)=\{(j,i):(i,j)\in E\}$, so $A=(A^{\mathrm{eco}})^{\top}$. We show that,
for the same parameters and random draws, the GPPM produces
$E^{\mathrm{eco}}$ exactly when our generator produces $R(E^{\mathrm{eco}})$,
with identical levels, trophic levels and $q$.

\begin{table}[h]
\centering
\small
\begin{tabular}{@{}p{0.15\linewidth}p{0.39\linewidth}p{0.39\linewidth}@{}}
\toprule
 & GPPM, ecological ($E^{\mathrm{eco}}$) & This paper ($E$) \\
\midrule
Tree phase & Newcomer $n$ picks prey $j_n$; adds $(j_n,n)$
           & Newcomer $n$ picks junior $j_n$; adds $(n,j_n)$ \\[2pt]
Excess edges & Basal nodes never predators (heads); weight uses
               $\hat s_{\mathrm{pred}}-\hat s_{\mathrm{prey}}-1$
             & Basal nodes never seniors (tails); weight uses
               $\ell_{\mathrm{src}}-\ell_{\mathrm{tgt}}-1$ \\[2pt]
Trophic level & Mean over in-neighbours, $\div\,k^{\mathrm{in}}$
              & Mean over out-neighbours, $\div\,k^{\mathrm{out}}$ \\
\bottomrule
\end{tabular}
\caption{The three direction reversals relative to the GPPM.}
\label{tab:reversals}
\end{table}

\emph{(i) Tree phase.} In the GPPM, newcomer $n$ picks $j_n$ uniformly from
existing nodes as its prey, adds $(j_n,n)$, and takes $\hat s_n=\hat
s_{j_n}+1$. We draw the same $j_n$, add $R((j_n,n))=(n,j_n)$, and assign the
same level. In both, the newcomer does the picking and sits one level above
its pick; only the drawn arrow differs. So $\ell\equiv\hat s$, and $D(G)$
agrees.

\emph{(ii) Excess edges.} A GPPM candidate $(i,j)$ (prey $i$, predator $j$)
needs $j$ non-basal. Under $R$ the predator is the source of $(j,i)$, so
``basal nodes are never predators'' becomes ``basal nodes are never sources'',
the condition $\ell_i\ge2$ in Eq.~\eqref{eq:candidates}. The GPPM weight
$\exp[-(\hat s_j-\hat s_i-1)^2/2T^2]$ of $(i,j)$ equals our $w_{ji}$. $R$ is
thus a weight-preserving bijection between candidate sets, and the sampling
distributions correspond edge for edge.
% CHECK against Klaise & Johnson (2016): (a) excess pairs drawn jointly with
% P_ij ∝ kernel (vs. predator first, then prey); (b) self-loops, tree edges and
% duplicates excluded, reciprocal pairs allowed. If either differs, move it to
% the departures sentence below.

\emph{(iii) Trophic levels.} The GPPM sets
$s_j=1+\frac{1}{k^{\mathrm{in}}_j}\sum_i A^{\mathrm{eco}}_{ij}s_i$ for
non-basal $j$. Substituting $A^{\mathrm{eco}}_{ij}=A_{ji}$ and
$k^{\mathrm{in}}_j=\sum_iA_{ji}=k_j$ gives Eq.~\eqref{eq:s} exactly, and the
basal sets coincide ($k^{\mathrm{in}}_j=0\iff k_j=0$). The trophic distance
$s_j-s_i$ of GPPM edge $(i,j)$ equals $x_{ji}$ (Eq.~\eqref{eq:x}) of the
reversed edge, so the multisets of distances, and hence $q$, are identical.
$\square$

The remaining departures from the GPPM are not reversals: $T$ is set per
realization from $\hat T$, only the largest component is kept, and excess
edges use successive sampling.
% TODO: numerical side-by-side check -- implement the GPPM in ecological
% orientation with the same RNG consumption, run with a shared seed, assert
% E == R(E_eco) and equal q over a grid of (N, T_hat). One sentence here.

\subsection{Normalized kernel width}
\label{sec:that}

\paragraph{Definition.}
For each realization, $D(G)=\max_i\ell_i$ is computed from the tree-phase
levels, after the tree phase and before any excess edge is drawn, and
$T=\hat T\,D(G)$. $D$ is a property of the tree, not of the trophic levels $s$
of the final graph. It is random, so $T$ is random at fixed $(N,\hat T)$.

\paragraph{Justification.}
The kernel acts on the integer gap $\Delta=\ell_i-\ell_j$. With
$\ell_i\in\{2,\dots,D\}$ and $\ell_j\in\{1,\dots,D\}$, the admissible gaps are
$\Delta\in\{2-D,\dots,D-1\}$, so $|\Delta-1|\le D-1$. In terms of the
normalized gap $u=\Delta/D\in(-1,1)$, Eq.~\eqref{eq:kernel} reads
\begin{equation}
  w=\exp\!\left[-\frac{(u-1/D)^2}{2\hat T^2}\right],
  \qquad
  \frac{w_{\min}}{w_{\max}}
  =\exp\!\left[-\frac{(1-1/D)^2}{2\hat T^2}\right]\ \xrightarrow{D\to\infty}\ e^{-1/(2\hat T^2)} .
  \label{eq:kernel-normalized}
\end{equation}
So $\hat T$ is the kernel width in units of the hierarchy's height. Over the
admissible gaps the kernel has the same shape at every $N$, up to an $O(1/D)$
shift of its peak, and the weight on the most extreme gap relative to the peak
is nearly independent of $N$ (about $0.61$ at $\hat T=1$, $0.14$ at
$\hat T=0.5$). With $T$ fixed, that ratio is $\exp[-(D-1)^2/2T^2]$, which falls
as $D$ grows with $N$ (mean $D\approx7$ at $N=10^2$, $\approx11$ at $N=10^4$).
Larger networks would then be more coherent by construction, confounding $N$
with the structural parameter whose effect on scaling we measure. We use the
realized $D(G)$ rather than a smooth function of $N$ (we also tried $\sqrt N$
and $\ln N$) because $D$ is what bounds $\Delta$ in a given network, and it
varies by several levels between realizations at the same $N$.

\paragraph{Limits.}
$\hat T$ fixes the kernel, not the realized gap distribution, which also
depends on how many candidates each level pair holds. Those counts come from
the tree's level occupancy, which is set by $b$: empirically $\ell-1$ is close
to Poisson with mean $\ln(1/b)$ at every $N$, so most nodes sit in the lowest
few levels and $D$ is set by a sparse top tail. We therefore measure $q$ on
every network rather than assume that fixing $\hat T$ fixes the incoherence.

\subsection{Trophic levels and incoherence}
\label{sec:coherence}

Trophic levels are computed on the final graph. With basal set
$\mathcal B=\{i:k_i=0\}$ and $F=V\setminus\mathcal B$, a node's trophic level
is one more than the mean level of the nodes it has authority over:
\begin{equation}
  s_i=1\quad(i\in\mathcal B),\qquad
  s_i=1+\frac{1}{k_i}\sum_jA_{ij}s_j\quad(i\in F).
  \label{eq:s}
\end{equation}
With $K_F=\operatorname{diag}(k_i)_{i\in F}$ and $r_i$ the number of basal
juniors of $i$, this is the linear system
$(K_F-A_{FF})\,s_F=k_F+r_F$. Every row of $K_F-A_{FF}$ is weakly diagonally
dominant, strictly so when $r_i>0$, and every $i\in F$ reaches such a row along
tree edges. The matrix is therefore weakly chained diagonally dominant and
nonsingular, so the solution is unique. We solve it with BiCGSTAB (relative
tolerance $10^{-10}$), with a sparse direct fallback.

Each edge has trophic distance $x_{ij}=s_i-s_j$. Because every edge has its
source in $F$, Eq.~\eqref{eq:s} gives
\begin{equation}
  \sum_{(i,j)\in E}x_{ij}=\sum_{i\in F}\Bigl(k_is_i-\sum_jA_{ij}s_j\Bigr)=\sum_{i\in F}k_i=|E|,
  \label{eq:x}
\end{equation}
so the mean distance is exactly one (numerically within $5\times10^{-11}$ in
every realization). The trophic incoherence parameter is the standard
deviation of the distances:
\begin{equation}
  q=\sqrt{\frac{1}{|E|}\sum_{(i,j)\in E}x_{ij}^2-1}\, .
  \label{eq:q}
\end{equation}
$q=0$ if and only if every edge spans exactly one trophic level, in which case
$s=\ell$.

\subsection{Consensus dynamics and scaling}
\label{sec:dynamics}

\paragraph{Influence weights.}
Every tie carries influence both ways; its direction only decides which side
is amplified. The weight agent $i$ places on agent $j$ is
\begin{equation}
  M_{ij}=[i=j]+[(i,j)\in E]+\alpha\,\bigl[(j,i)\in E\wedge(i,j)\notin E\bigr],
  \qquad W=\operatorname{diag}(M\mathbf1)^{-1}M .
  \label{eq:W}
\end{equation}
On an unreciprocated tie the senior weights the junior at $1$ and the junior
weights the senior at $\alpha$. Reciprocated ties and self-ties have weight
$1$.

\paragraph{Dynamics and consensus time.}
Opinions follow DeGroot updating~\cite{degroot1974reaching}, $x(t+1)=Wx(t)$.
The support of $M$ is the connected undirected projection of $G$ plus
self-loops, so $W$ is primitive: $\lambda_1=1$ is simple, all other
eigenvalues have modulus below one, and opinions reach consensus. For
$\alpha\ne1$, $W$ is generally not reversible: the Kolmogorov product around a
cycle is $\alpha^{m_\uparrow-m_\downarrow}$, where $m_\uparrow$ and
$m_\downarrow$ count the unreciprocated ties traversed junior-to-senior and
senior-to-junior. Eigenvalues may therefore be complex and are ordered by
modulus. Disagreement decays asymptotically as $|\lambda_2|^t$, and we define
consensus time as the steps needed for it to shrink by a factor $\delta$:
\begin{equation}
  \tau=\frac{\ln(1/\delta)}{\ln(1/|\lambda_2|)},\qquad\delta=10^{-6}.
  \label{eq:tau}
\end{equation}
$|\lambda_2|$ is computed with ARPACK, with a dense fallback that was never
needed. We do not iterate the dynamics: $\tau$ is a spectral estimate and
omits the prefactor set by $x(0)$ and the conditioning of $W$'s eigenvectors.
% TODO: validate tau against simulated iterate counts (stop when
% max_i x_i - min_i x_i < delta) on a subsample of (N, T_hat) cells.

\paragraph{Scaling exponent.}
We sweep 20 sizes log-spaced over $N\in[10^2,10^4]$ (the set $\mathcal N$) and
$\hat T\in\{0.1,0.2,\dots,1.0\}$, with $b=0.1$, $c\approx7.12$, $\alpha=2$ and
$M=20$ independent realizations per $(N,\hat T)$ cell, for 4{,}000 networks.
All randomness comes from one PCG64 stream (seed 21) consumed in sweep order.
For each cell we take the median consensus time
$\tilde\tau(N,\hat T)=\operatorname{median}_m\tau^{(m)}$ and fit, by ordinary
least squares over $\mathcal N$,
\begin{equation}
  \ln\tilde\tau(N,\hat T)=\ln C(\hat T)+\gamma(\hat T)\ln N ,
  \label{eq:gamma}
\end{equation}
so $\tilde\tau\propto N^{\gamma}$, and $\gamma\approx0$ means consensus time
does not grow with group size. Each $\hat T$ is summarized by its mean
incoherence over all sizes and realizations,
$\bar q(\hat T)=(|\mathcal N|M)^{-1}\sum_{N}\sum_{m}q^{(m)}(N,\hat T)$, and we
report $\gamma$ against $\bar q$.
% TODO: q varies with N at fixed T_hat, so q-bar is an ensemble average.
%       Consider also reporting q at N_max, or gamma against q within each N.
% TODO: no uncertainty on gamma yet -- bootstrap over realizations per cell.

% ---------------------------------------------------------------------
% BibTeX (move to your .bib):
%
% @article{johnson2014trophic,
%   author  = {Johnson, Samuel and Dom{\'\i}nguez-Garc{\'\i}a, Virginia and
%              Donetti, Luca and Mu{\~n}oz, Miguel A.},
%   title   = {Trophic coherence determines food-web stability},
%   journal = {Proceedings of the National Academy of Sciences},
%   volume  = {111}, number = {50}, pages = {17923--17928}, year = {2014}}
%
% @article{klaise2016neurons,
%   author  = {Klaise, Janis and Johnson, Samuel},
%   title   = {From neurons to epidemics: How trophic coherence affects
%              spreading processes},
%   journal = {Chaos}, volume = {26}, number = {6}, pages = {065310},
%   year    = {2016}}
%
% @article{degroot1974reaching,
%   author  = {DeGroot, Morris H.},
%   title   = {Reaching a consensus},
%   journal = {Journal of the American Statistical Association},
%   volume  = {69}, number = {345}, pages = {118--121}, year = {1974}}
% ---------------------------------------------------------------------
